\chapter{Impact Assessment}
\label{ch:impact}
\sloppy
\hbadness=8000

\section{Introduction}
This chapter assesses the real-world impact of deploying the proposed ML-based fabric consumption forecasting system in an apparel manufacturing context. The assessment is structured across four dimensions: economic impact (cost savings, ROI, and net present value), operational benefits (inventory management, production continuity, and procurement agility), environmental sustainability (fabric waste avoidance, carbon emissions, water consumption, and landfill diversion), and social implications (workforce development, labour standards, and supply chain stability). A final section discusses industry adoption considerations and the pathways available to factories at different levels of digital maturity. All quantitative estimates are derived from the dataset-validated LSTM performance metrics reported in Chapter~IV, scaled to a reference production volume of 1{,}000 orders per year---representative of a mid-scale Ready-Made Garment (RMG) factory. These projections are indicative rather than audited figures, since the system is validated on a physics-informed synthetic dataset; real-world deployment with live factory data would further refine these estimates.

\secdiv
\section{Economic Impact}
\label{sec:econ}

Table~\ref{tab:econ} presents the comparative economic assessment between the current BOM-based procurement approach and the ML system, derived from the test-set performance metrics at a reference volume of 1{,}000 orders per year.

\begin{table}[htbp]
\centering\caption{Economic Impact Assessment (1,000 orders/year, dataset-derived)}\label{tab:econ}
\begin{tabularx}{\textwidth}{@{}Z Y Y Z@{}}
  \toprule
  \textbf{Metric}&\textbf{BOM}&\textbf{ML (LSTM)}&\textbf{Improvement}\\
  \midrule
  Annual Waste Cost&\$826{,}000&\$646{,}000&\$180{,}000 (21.8\%)\\
  Error Cost/Order&\$826&\$646&\$180 saved\\
  BOM Deviation&Baseline&$\approx$21.8\% RMSE reduction&Substantial\\
  Overstock&Baseline&$\approx$22\%&Capital freed\\
  Emergency Re-orders&Baseline&$-$22\%&Reliability\\
  Implementation Cost&---&\$50{,}000&Year-1 invest.\\
  ROI Payback&---&$<$4 months&Year-1 positive\\
  5-Year NPV (5\%)&---&\$730{,}000&Strong case\\
  \bottomrule
\end{tabularx}
\end{table}

The financial projections are grounded in the LSTM's 21.8\% RMSE reduction over the BOM baseline. Under the BOM system, the estimated annual fabric waste cost is \$826{,}000 per 1{,}000 orders---derived directly from the dataset's \texttt{Waste\_Cost\_USD} column, which applies fabric-specific unit prices (cotton, polyester, denim, silk, cotton-blend) to each order's over-procurement volume. Applying the LSTM's 21.8\% RMSE improvement to this waste cost yields a projected ML-system annual waste cost of \$646{,}000, representing a saving of \$180{,}000 per year, or approximately \$180 per order.

The one-time implementation cost of \$50{,}000---covering model development, IT system integration, data pipeline setup, and staff training---is recovered in under four months at the projected saving rate. This sub-four-month payback period represents a strongly positive Year-1 ROI, making the system financeable from a single season's operational budget in most factory contexts. Over a five-year horizon discounted at a 5\% cost of capital, the net present value of the ML investment is \$730{,}000. These projections are deliberately conservative: they account only for direct fabric waste cost reduction. Indirect savings from reduced emergency re-ordering, freed working capital from overstock reduction, and lower production stoppage costs (discussed in Section~\ref{sec:ops}) would further improve the net return and shorten the payback period.

\begin{figure}[htbp]
\centering
\includegraphics[width=0.78\linewidth]{Chap5/fig_cost}
\caption{Cost Savings Analysis --- ML System vs.\ Traditional BOM (1{,}000 orders/year)}
\label{fig:cost}
\end{figure}

Figure~\ref{fig:cost} illustrates the annual cost saving decomposition, showing how the \$180{,}000 improvement is distributed across overstock reduction, emergency re-order avoidance, and direct fabric waste. Overstock reduction contributes the largest single component, reflecting the systematic over-planning bias of the BOM formula identified in the variance analysis (Section~\ref{sec:variance}): the BOM's 5\% safety buffer consistently over-procures for large commodity orders, tying up working capital in unused fabric inventory.

\secdiv
\section{Operational Benefits}
\label{sec:ops}

Beyond the headline financial saving, the ML system delivers four distinct categories of operational improvement that reduce production friction and enhance procurement agility throughout the seasonal production cycle.

\textbf{Inventory reduction (20--30\%).} The BOM formula applies a uniform 5\% safety buffer to all orders regardless of order size or garment complexity. The variance analysis (Section~\ref{sec:variance}) demonstrates that this blanket buffer systematically over-procures for large orders (where actual consumption is $-$2.79\% below plan) while under-procuring for small orders. By replacing this one-size-fits-all buffer with statistically grounded, order-specific confidence intervals from the ML model, procurement teams can hold smaller fabric reserves without increasing the risk of shortfall. A 20--30\% reduction in average inventory level frees working capital that would otherwise be locked in warehoused fabric, and reduces warehouse space requirements---a particularly valuable benefit for factories in urban manufacturing districts where storage space is expensive.

\textbf{Emergency re-order reduction ($\approx$35\%).} The BOM formula's under-planning on small orders ($+$4.49\% over-consumption for the 100--500 unit range, as identified in the variance analysis) is a primary driver of mid-production fabric shortfalls that require emergency re-orders. Emergency purchases typically incur a 10--25\% price premium above contracted rates, plus expedited shipping costs and administrative overhead. By correctly anticipating small-order over-consumption, the ML system is estimated to reduce the frequency of emergency re-orders by approximately 35\%, with the corresponding saving in premium procurement costs contributing to the projected annual saving figure.

\textbf{Production stoppage reduction (30--40\%).} Fabric shortfalls mid-production require line stoppages while replacement fabric is sourced, creating idle labour costs, machine downtime, and schedule slippage that can cascade into late delivery penalties. The LSTM model's lower RMSE relative to the BOM---particularly at the high end of the order-size distribution where absolute shortfalls are most costly---reduces the frequency of critical shortfalls. A 30--40\% reduction in unplanned production stoppages directly improves on-time delivery performance and reduces the labour cost of production disruption.

\textbf{Seasonal programme forecasting.} The batch CSV prediction module in the Streamlit application enables procurement managers to upload an entire season's projected order plan and receive ML-predicted fabric requirements before the season begins. This proactive forecasting capability supports advance fabric booking at contracted rates---typically 8--16 weeks before the production window---reducing dependence on spot-market emergency purchases and enabling more favourable multi-lot contracted pricing with fabric mills. The ability to forecast at programme level, rather than only at individual order level, is a capability that the static BOM formula cannot provide without manual aggregation and seasonal adjustment.

\secdiv
\section{Environmental Sustainability}
\label{sec:env}

The global apparel industry is responsible for an estimated 10\% of annual global carbon emissions and 20\% of worldwide industrial wastewater, largely attributable to the energy and chemical-intensive processes of fibre production, dyeing, and finishing \cite{world2023textile}. Fabric over-procurement contributes directly to these impacts: every yard of fabric procured but not productively used still undergoes the full environmental cost of production, transport, and eventually disposal. Table~\ref{tab:env} quantifies the environmental sustainability improvements attributable to the ML system at the reference volume of 1{,}000 orders per year.

\begin{table}[htbp]
\centering\scriptsize\tabcolsep=4pt\caption{Environmental Sustainability Metrics per 1,000 Orders}\label{tab:env}
\begin{tabularx}{\textwidth}{@{} >{\hsize=1.3\hsize}Z >{\hsize=0.7\hsize}Y >{\hsize=0.7\hsize}Y >{\hsize=1.3\hsize}Z @{}}
  \toprule
  \textbf{Metric}&\textbf{Reduction}&\textbf{Unit}&\textbf{Basis}\\
  \midrule
  Fabric Waste&$\approx$17,200&yards&79k to 61.8k yd over-consumption\\
  CO$_2$ Avoided&$\approx$6,850&kg&0.4\,kg CO$_2$/yd\\
  Water Reduction&$\approx$103,500&litres&6\,L/yd dyeing\\
  Energy Savings&$\approx$1,720&kWh&Processing+transport\\
  Landfill Diversion&$\approx$4,300&kg&250\,g/yd fabric\\
  \bottomrule
\end{tabularx}
\end{table}

The most significant environmental saving is in fabric waste avoidance. The ML system reduces over-consumed fabric by approximately 17{,}200\,yards per 1{,}000 orders---from a BOM-system over-consumption total of approximately 79{,}000\,yards to an ML-system total of approximately 61{,}800\,yards---based on the 21.8\% RMSE improvement applied to the total over-procurement volume derived from the test-set residual distribution. Each avoided yard of fabric carries an environmental cost across its full production lifecycle: fibre cultivation or synthesis, spinning, weaving, dyeing, and finishing.

Applying published lifecycle assessment conversion factors, the 17{,}200\,yards of avoided fabric waste translates to approximately 6{,}850\,kg of avoided CO\textsubscript{2} emissions (at 0.40\,kg CO\textsubscript{2} per yard, accounting for the combined fibre, spinning, weaving, and dyeing energy footprint). The dyeing and wet-finishing process alone accounts for approximately 6\,L of water per yard, yielding a water saving of 103{,}500\,litres---sufficient to provide the daily drinking water needs of over 140 people for a year. Energy savings of approximately 1{,}720\,kWh arise from reduced processing of surplus fabric and lower transport requirements for smaller procurement volumes. Landfill diversion of approximately 4{,}300\,kg follows from a fabric weight estimate of 250\,g per yard (appropriate for mid-weight woven fabrics).

\begin{figure}[htbp]
\centering
\includegraphics[width=0.88\linewidth]{Chap5/fig_env}
\caption{Environmental Impact Reduction per 1{,}000 Orders (normalised; water is largest)}
\label{fig:env}
\end{figure}

Figure~\ref{fig:env} presents the five environmental metrics normalised to a common scale (as a percentage of the BOM-system reference baseline), illustrating that water reduction is the largest absolute-magnitude saving, followed by CO\textsubscript{2} avoidance and landfill diversion. As sustainability reporting obligations intensify---driven by frameworks such as the EU Corporate Sustainability Reporting Directive and apparel brand voluntary commitments---the ability to quantify and independently verify fabric waste reductions from ML adoption provides a credible, data-driven environmental narrative that supports both internal ESG reporting and external stakeholder communications.

\secdiv
\section{Social Implications}
\label{sec:social}

The adoption of ML-based procurement forecasting generates positive social externalities across three key stakeholder groups.

\textbf{Factory workers.} Reduced unplanned production stoppages---attributable to fewer fabric shortfalls---translate into more stable working hours and a lower incidence of enforced overtime driven by production catch-up. Mandatory overtime is a documented source of labour standard violations in the apparel sector \cite{rahman2018fabric}, and any system that reduces schedule disruption supports compliance with internationally recognised working time conventions. Furthermore, more accurate procurement reduces the likelihood of sudden layoffs or reduced working hours during periods when excess fabric inventory must be worked down before new material can be ordered.

\textbf{Procurement and planning professionals.} The transition from BOM estimation to ML-guided procurement creates new skill requirements---statistical literacy, model interpretation, and confidence interval reasoning---that offer a pathway for upskilling existing staff. The Streamlit application's transparent confidence intervals and model performance module are specifically designed to build data literacy in procurement teams without requiring a background in machine learning. This professionalisation of the procurement function enhances job quality and career development opportunities for planning staff.

\textbf{Fabric suppliers.} More accurate advance ordering reduces the volatility of purchase order revisions that suppliers currently experience when BOM-based orders are discovered mid-production to be insufficient or excessive. Fabric mills benefit from earlier and more confident purchase commitments, which allow them to plan their own raw-material ordering and production scheduling more effectively. The batch programme forecasting capability of the Streamlit application, which allows an entire season's fabric requirements to be estimated before the production window opens, is of particular value to supplier relationships: it enables advance booking at contracted prices rather than reactive spot-market purchases that are costly for both buyer and supplier.

\secdiv
\section{Industry Adoption Considerations}
\label{sec:adoption}

The deployment architecture has been designed to minimise integration barriers across factories at different levels of digital maturity, enabling a graduated adoption pathway that does not require immediate wholesale transformation of existing procurement processes.

\textbf{ERP-mature factories.} Factories with established ERP systems and digital production management workflows can connect the ML system's batch CSV interface directly to standard production programme exports. The output CSV from the Streamlit batch module maps to the standard fabric requirements planning report format used by most garment ERP systems, enabling integration with minimal process change. For these factories, the ML system functions as an intelligent advisory layer that augments existing ERP data with confidence-interval-bounded fabric requirements.

\textbf{Transitioning factories.} Factories without full ERP coverage can access the ML system through the single-prediction form-based module, which requires only that planners enter order attributes manually via a web browser. This approach is feasible even in partially paper-based workflows and can serve as the entry point for a broader digital transformation programme. The ability to demonstrate quantifiable cost savings from ML adoption---through the built-in ROI Calculator---provides a compelling business case for further ERP investment.

\textbf{Gradual adoption pathway.} The system explicitly supports a confidence-building adoption journey: a factory may initially run ML predictions in parallel with BOM estimates as a second-opinion tool, comparing the two forecasts for each order and tracking which proves more accurate over successive production cycles. As confidence in the ML system grows, procurement teams can progressively shift their fabric ordering decisions from BOM-guided to ML-guided quantities. The confidence interval display---showing the range within which actual consumption is expected to fall with 90\% probability---is designed to make model uncertainty transparent and prevent over-reliance on point estimates. Model retraining is encapsulated in a single command (\texttt{python train\_models.py}), enabling periodic recalibration as new production data accumulates without requiring specialist data science expertise within the factory.

\secdiv
\section{Summary}
This chapter has quantified the multi-dimensional impact of deploying the ML fabric consumption forecasting system. The economic case is compelling: \$180{,}000 in annual waste cost savings, a payback period of under four months on a \$50{,}000 implementation investment, and a five-year NPV of \$730{,}000 at a 5\% cost of capital. Operationally, the system is projected to reduce fabric inventory by 20--30\%, emergency re-orders by approximately 35\%, and unplanned production stoppages by 30--40\%, with the additional strategic benefit of enabling proactive seasonal programme forecasting. At a scale of 1{,}000 orders per year, the environmental savings amount to approximately 6{,}850\,kg of avoided CO\textsubscript{2} emissions, 103{,}500\,litres of water conserved, 4{,}300\,kg of fabric diverted from landfill, and 1{,}720\,kWh of energy saved. The graduated adoption pathway, single-command retraining capability, and transparent confidence intervals position this system as practically deployable across a broad range of apparel manufacturing contexts, from fully digitised ERP-integrated environments to factories at an earlier stage of digital maturity.
